% ============================================================================
% INTEGRATED EVALUATION (core) — for read-through before merging into main.tex.
% Locked config: RESISC45 n=1300 all-Q4; P_cap=15W shared-budget abstraction;
% P_base=7W nominal (lit. 3-10W) + swept; MAX peak calibration; V NOT pinned.
% Numbers aligned across all subsections. NOTE: sensitivity/ISL/stability
% subsections still carry OLD-config numbers and must be re-run separately (flagged).
% ============================================================================

\subsection{Hardware Calibration and Model Selection}\label{EV_CALIB}
We profile seven vision-language models on RESISC45 (1300 images, 19 to 41 per
category), each at 4-bit quantization on an NVIDIA Jetson Orin NX in its 15\,W
power mode. Latency is timed per image, power is sampled at 50\,ms with
\texttt{tegrastats}, and energy integrates power over the inference window.
Fig.~\ref{fig:tradeoff} places every model on the accuracy-energy plane. Four
models lie on the Pareto frontier: Qwen2-VL-2B (0.394 accuracy, 3.93\,J),
Qwen2.5-VL-3B (0.456, 5.72\,J), Qwen2.5-VL-7B (0.549, 9.30\,J), and InternVL3-8B,
which reaches the highest accuracy (0.562) at the highest energy (15.00\,J). The
other three are dominated: SmolVLM2-2.2B (0.310), InternVL2.5-4B (0.548), and
gemma-3-4b (0.538, 37.79\,J).

We adopt the Qwen family as the deployment backbone for an architectural reason,
not a per-point accuracy judgement. Qwen2-VL-2B, Qwen2.5-VL-3B, and
Qwen2.5-VL-7B share one architecture and one activation interface, so they form a
single three-tier ladder that maps onto the deployment tiers of this work: a
light single-satellite model, a heavier single-satellite model, and a split-only
model. Their uniform transformer-block structure keeps the activation dimension
constant at any partition point, which is the precondition for the two-stage
split. InternVL3-8B is an isolated 8B model of a different architecture: it does
not compose with the on-board 2B and 3B tiers into one routable family, its
larger size is even less feasible on a single satellite, and splitting it is
heavier; its accuracy lead therefore does not change the deployment choice.
InternVL2.5-4B is within statistical noise of Qwen2.5-VL-7B (0.548 versus 0.549).
SmolVLM2-2.2B is the weakest model and is dominated by the smaller Qwen2-VL-2B,
which shows that a model's frontier position is set by architecture and training
rather than by parameter count alone.

The split pipeline is required by a hardware constraint that does not depend on
any accuracy comparison. A 4-bit Qwen2.5-VL-7B model, with its vision encoder and
key-value cache, does not fit in the memory left on an 8\,GB satellite after the
on-board sensing pipeline, so the two-satellite split is the only way to run the
7B tier at all.

\subsection{Peak-Power Calibration}\label{EV_PEAK}
The instantaneous power cap is the second structural constraint. We model it as a
shared power budget $P^{\text{cap}}{=}15$\,W, taking the Jetson 15\,W envelope as
the nominal value; the platform baseline $P_{\text{base}}$ and the inference peak
draw from this budget. This is an abstraction: a real satellite separates the
platform power system from the compute payload, and a higher-fidelity power model
is left to future work. We calibrate the peak bound $\hat{P}_{im}^{\text{CP}}$ as
the maximum per-inference peak observed over the benchmark, a conservative upper
bound appropriate for a hard constraint that must never be violated. Among the
deployed configurations only the 7B tier shows a sharp transient: its peak reaches
6.37\,W against a 3.77\,W median (a $1.9\times$ ratio), traced to a sustained
compute-bound prefill phase, whereas the 2B and 3B peaks (3.75 and 4.02\,W) sit
close to their own high quantiles. The 7B peak plus a few-watt platform baseline
approaches the 15\,W budget, so a scheduler that commits the split without
checking the peak can drive the combined draw past the cap. This is the empirical
basis for the per-slot peak shield.

\subsection{Simulation Setup}\label{EV_SETUP}
Each slot is $\tau{=}1$\,s; one orbit is 5700 slots (60\% sunlit, 40\% eclipse) and
we run three orbits. The constellation has four computing satellites and six
sensing sources with heterogeneous, data-driven quality profiles taken from the
per-category RESISC45 measurements. The platform baseline continuously draws
$P_{\text{base}}$ from the battery; in eclipse it is borne entirely by the battery
and forms the bulk of the eclipse energy budget. We take $P_{\text{base}}{=}7$\,W
as a nominal value in the reported CubeSat range (3--10\,W for attitude control and
housekeeping) and sweep it in Section~\ref{EV_SENS}; the solar panel
($P^{\text{solar}}{=}13$\,W) yields an orbit-average harvest just above the
platform draw, leaving a thin margin for paced inference, and the battery
(18\,kJ) is sized to carry the platform through eclipse. All policies draw from
one shared feasible set (split included, the ISL handoff check, the per-slot
energy check, and the peak check), so they differ only in the selection rule;
quality ties are broken uniformly at random.

\subsection{Constrained-Fair Comparison}\label{EV_MULTI}
Fig.~\ref{fig:frontier} and Table~\ref{tab_split} report the four-satellite
comparison on the quality-fairness plane (10-seed means). The proposed scheduler
is not a single point but a curve swept by the Lyapunov weight $V$: small $V$
favours the max-min objective (min-fill up to 0.175), large $V$ favours total
throughput (up to 3850), and we pin no single $V$. The baselines are fixed points.
The reading is threefold. First, the proposed scheduler is the only policy that is
simultaneously deployable and able to run the 7B split: it holds downtime at 0.4\%
across the whole $V$ range while running 179 to 206 splits, whereas the two other
deployable policies (Greedy-E and Static) never split. Second, at small $V$ it
exceeds the fair baseline Greedy-E on min-fill (0.175 versus 0.156, a $5\sigma$
gap) at comparable total quality (3254 versus 3193, a difference within
$1.6\sigma$, not significant). Third, the policies that reach both high fairness and high throughput
(Greedy-Q at min-fill 0.398, Random at 0.354) do so only by blacking out (7.0\% and
5.3\% downtime), which violates the uptime service-level agreement (SLA); the
deployability partition is invariant for any SLA threshold in the wide band
$(0.8\%, 5.3\%)$. No policy reaches high fairness and high throughput while
remaining deployable, which is the quality-fairness conflict made visible as the
empty upper region of Fig.~\ref{fig:frontier}. Static reaches higher raw
throughput (5102) at low fairness (0.119); the proposed scheduler trades this
throughput for the ability to navigate to the fair region and to run the top tier.

\begin{table}[!t]
\centering\footnotesize
\setlength{\tabcolsep}{3pt}
\caption{\label{tab_split}Four-satellite comparison (three orbits, 10-seed
mean$\,\pm\,$std). All policies draw from the same shared feasible set. The
proposed scheduler is reported at three values of the Lyapunov weight $V$ to show
the fairness--throughput knob; no single $V$ is selected (Fig.~\ref{fig:frontier}).
Blackout above the uptime SLA marks a policy undeployable.}
\begin{tabular}{lcccccc}
\toprule
Policy & \makecell{Peak\\viol.} & \makecell{Down\\time} & \makecell{Serv.\\rate} & \makecell{7B\\splits} & \makecell{Min-\\fill} & \makecell{Total\\qual.} \\
\midrule
Proposed ($V{=}1$)   & 0 & 0.4\% & 0.525 & $206$ & $\mathbf{0.175}$ & $3254$ \\
Proposed ($V{=}10$)  & 0 & 0.4\% & 0.526 & $201$ & $0.169$ & $3311$ \\
Proposed ($V{=}100$) & 0 & 0.4\% & 0.534 & $179$ & $0.124$ & $3636$ \\
\midrule
Greedy-Q & 0 & \textbf{7.0\%} & 0.765 & 5299 & 0.398 & 6083 \\
Greedy-E & 0 & 0.8\% & 0.780 & 0 & 0.156 & 3193 \\
Random   & 0 & \textbf{5.3\%} & 0.778 & 3333 & 0.354 & 5022 \\
Static   & 0 & 0.8\% & 0.780 & 0 & 0.119 & 5102 \\
\bottomrule
\end{tabular}
\end{table}

\subsection{Capability and Routing}\label{EV_ROUTE}
The heterogeneity of the sources makes model choice a routing problem rather than a
size ranking. To avoid selecting categories after seeing the data, we fixed a test
in advance: a category requires the split-only 7B model if 3B scores at most 0.35,
7B scores at least 0.70, and the difference is significant under a two-sided Fisher
exact test after Benjamini-Hochberg correction across all 45 categories. Five
categories pass. Four of them, however, are non-monotone in model size: the smaller
2B model scores 0.79 to 0.89, 3B collapses to 0.04 to 0.34, and 7B recovers to 0.78
to 0.83. This is a category-specific weakness of the 3B model, and because 2B also
runs on a single satellite the best single-satellite option already serves these
four. Under the stricter criterion that both single-satellite models fail
($\max(\text{2B},\text{3B})\le 0.35$) while 7B reaches 0.70, which we mark as a
deviation from the pre-registration, only one category (mountain) qualifies, and it
passes Benjamini-Hochberg but not Bonferroni. The aggregate accuracy gain of 7B
over 3B is only 0.093, and the per-category sample is small (19 to 41 images,
$95\%$ interval about $\pm0.18$), so we rest the case for the split on the memory
constraint rather than on any single category. The value of the data is that model
size alone is a poor guide to configuration: most categories are served by a
single-satellite model, several on which 3B fails are recovered by the smaller 2B,
and the scheduler routes on the measured per-source quality rather than on size.

\subsection{The Fairness--Throughput Conflict}\label{EV_VCONF}
Fig.~\ref{fig:vconf} traces the two objectives against $V$. Total delivered quality
rises with $V$, so the $O(1/V)$ improvement of Theorem~\ref{thm_gap} holds in the
multi-satellite setting. Min-fill, in contrast, falls with $V$. This reverses its
single-satellite behaviour, where min-quality rises with $V$ because there is no
fairness conflict; here the heterogeneous sources create one. The reversal does not
violate Theorem~\ref{thm_gap}, which bounds total quality rather than min-fill, and
it makes explicit that $O(1/V)$ is a total-quality guarantee. The reversal is the
$V$-space signature of the quality-fairness conflict.

% ---- FLAGGED: to re-run with the locked config before final integration ----
% \subsection{Sensitivity Analysis} (P_base sweep = the 7B feasibility boundary
%   at ~8.63W; panel/battery/nsat sweeps) -- OLD-config numbers, MUST re-run.
% \subsection{Dynamic ISL Topology} -- OLD-config, re-run.
% \subsection{Stability Stress Tests} -- OLD-config, re-run.
